%%%%%%%%%%%%%%%%%%%%%%%%%%%%%%%%%%%%%%%%%%%%%%%%%%%%%%%%%%%%%%%%%%%%%%%%%%%%%%%%%%
\begin{frame}[fragile]\frametitle{}
\begin{center}
{\Large ओशो यांच्या शिकवणी}
\end{center}

{\tiny (Ref: ओशो --- \textit{व्यस्त जीवनातील अध्यात्म} (ईश्वर की खोज, मराठी अनुवाद))}

\end{frame}

%%%%%%%%%%%%%%%%%%%%%%%%%%%%%%%%%%%%%%%%%%%%%%%%%%%%%%%%%%%
\begin{frame}[fragile]\frametitle{३ स्तर: शरीर, मन, आत्मा}

	\begin{itemize}
	\item ३ स्तर : शरीर, मन, आत्मा
	\item शरीर: कष्ट असतात किंवा नसतात --- येथे सुख नसते
	\item मन: सुख असते किंवा नसते --- येथे कष्ट नसतात
	\item आत्मा: दु:ख नसल्यास उरतो तो फक्त आनंद
	\item हे पंचकोशच आहेत ---
		\begin{itemize}
		\item अन्नमय (शरीर), प्राणमय (श्वास --- एक सेतू)
		\item मनोमय (मन), विज्ञानमय (उच्च मन), आनंदमय (आत्मा)
		\end{itemize}
	\end{itemize}

{\tiny (Ref: ओशो --- \textit{व्यस्त जीवनातील अध्यात्म})}

\end{frame}

%%%%%%%%%%%%%%%%%%%%%%%%%%%%%%%%%%%%%%%%%%%%%%%%%%%%%%%%%%%
\begin{frame}[fragile]\frametitle{शरीर आणि मन}

	\begin{itemize}
	\item शरीराला एकतानता लागते --- सुरक्षिततेसाठी नवनवीन नको वाटते
	\item उदा: शरीराला दर दिवशी वेगवेगळ्या ठिकाणी झोप लागेल का?
	\item मनाला सतत नित्यनुतन लागत असते --- सुख क्षणभंगुर असते
	\item उदा: तोच तो शर्ट सारखा घालता का? नेहमी काहीतरी नवीन घालावेसे वाटते
	\end{itemize}

{\tiny (Ref: ओशो --- \textit{व्यस्त जीवनातील अध्यात्म})}

\end{frame}

%%%%%%%%%%%%%%%%%%%%%%%%%%%%%%%%%%%%%%%%%%%%%%%%%%%%%%%%%%%
\begin{frame}[fragile]\frametitle{आत्मा: सच्चिदानंद}

	\begin{itemize}
	\item एकदा सच्चिदानंद मिळाला की तो कायमचा राहतो
	\item आत्मा ते परमात्मा --- हे जीवनाचे ध्येय आहे
	\item त्यासाठी वेगळा वेळ काढण्याची अथवा काही विशेष कार्य करण्याची गरज नाही
	\item आपलेच काम अजगतेने, मन लावून उत्तमोत्तम करण्यानेच तो आनंद मिळायला लागतो
	\end{itemize}

{\tiny (Ref: ओशो --- \textit{व्यस्त जीवनातील अध्यात्म})}

\end{frame}
